\documentclass{resume} % Use the custom resume.cls style

\usepackage[left=0.4 in,top=0.4in,right=0.4 in,bottom=0.4in]{geometry} % Document margins
\newcommand{\tab}[1]{\hspace{.2667\textwidth}\rlap{#1}}
\newcommand{\itab}[1]{\hspace{0em}\rlap{#1}}
\name{Giacomo Fantoni} % Your name
% You can merge both of these into a single line, if you do not have a website.
\address{+39(328) 118-2628 \\ Localit\`a Negrano 6, 38123, Trento, Italy \\ \href{mailto:giacomo.fantoni@unitn.it}{giacomo.fantoni@unitn.it}}

\begin{document}

\begin{rSection}{}
   Dottorando in Scienze biomolecolari con focus sull'analisi di dati biomedici complessi, con focus sulle omiche
   Proponente del progetto JoinTheLab, con l’obiettivo di integrare l'esperienza tecnica nell'analisi di dati NGS complessi a una visione imprenditoriale mirata alla digitalizzazione dei processi di laboratorio.
   Impegnato nel definire nuovi standard di riproducibilità per il settore biotech attraverso il trasferimento tecnologico della ricerca accademica
\end{rSection}
%----------------------------------------------------------------------------------------
%	OBJECTIVE
%----------------------------------------------------------------------------------------

% \begin{rSection}{OBJECTIVE}

% {Software Engineer with 2+ years of experience in XXX, seeking full-time XXX roles.}


% \end{rSection}
%----------------------------------------------------------------------------------------
%	EDUCATION SECTION
%----------------------------------------------------------------------------------------

\begin{rSection}{Educazione}


{\bf Universit\`a degli studi di Trento} \hfill {Trento, Italia}\\
Corso di Dottorato in Scienze Biomolecolari \hfill {Nov. 2024 - In Corso}
\begin{itemize}
   \item Implementazione di un innovativo algoritmo di regressione basato sui ranghi per la predizione di variabili cliniche continue a partire da dati omici di varia natura.
   \item Processamento e analisi di dati omici, ottenendo risultati rilevanti in ambito biologico grazie all'applicazione di metodi di Machine Learning "Explainable".
   \item Presentazione dei risultati ottenuti a conferenze quali \href{https://www.bioinformatics.polimi.it/CIBB2025/}{CIBB 2025} e \ecvhighlight{\href{https://www.ds4hb.polimi.it/}{DS4HB}}.
\end{itemize}

{\bf Universit\`a degli studi di Trento} \hfill {Trento, Italia}\\
Laurea Magistrale in Biologia Quantitativa e Computazionale \hfill {Set. 2021 - Ott. 2023}
Tesi: "Computational modelling and simulation of the antennal lobe of the Honey Bee to uncover neuronal correlates of sleepAComputational modelling and simulation of the antennal lobe of the Honey Bee to uncover neuronal correlates of sleep". Relatore Prof. A. Haase. Voto Finale 110/110 con lode.
\begin{itemize}
   \item Corsi Rilevanti: Basi di Dati, Algoritmi e Strutture Dati, Ingegneria del software.
\end{itemize}

{\bf Universit\`a degli studi di Trento} \hfill {Trento, Italia}\\
Laurea Triennale in Informatica\hfill {Set. 2018 - Set. 2021}\\
Tesi: "Analisi di dati di trascrittomica RNA-seq da frazioni polisomiali e totali da una linea cellulare di cancro epiteliale. Relatore Prof. A. Inga. Voto Finale 108/110.
\begin{itemize}
   \item Corsi Rilevanti: Bioinformatics, Data mining, Genomics.
\end{itemize}


\end{rSection}


\begin{rSection}{Esperienza}

\textbf{Assegno di ricerca - Universit\`a degli studi di Trento, Dipartimento CIBIO} \hfill Novembre 2023 - Novembre 2024\\
PREDICTive Cardio-metabolic transcriptOMIC trajectorieS in the Barilla Offspring Follow-up STUDY: The PREDICT-OMICS Study \hfill \textit{Trento, Italia}
 \begin{itemize}
    \itemsep -3pt {}
     \item Implementazione di una pipeline automatica per il processamento di dati microarray provenienti dal Barilla Offspring Study.
     \item I risultati sono stati presentati come poster alla conferenza \href{https://bioinformatics.it/bits2024/}{BITS 2024}.
 \end{itemize}

\textbf{Tirocinio - Universit\`a degli studi di Trento, Dipartimento CIMEC} \hfill Marzo 2023 - Settembre 2023\\
Computational modelling and simulation of the antennal lobe of the Honey Bee to uncover neuronal correlates of sleep \href{https://github.com/giacThePhantom/genn-network-model}{github.com/giacThePhantom/genn-network-model} \hfill \textit{Rovereto, Italia}
 \begin{itemize}
    \itemsep -3pt {}
     \item Implementazione di un modello computazionale accelerato da GPU in python.
     \item Implementazione di una pipeline di machine learning per l'analisi dei dati sperimentali e del modello.
 \end{itemize}

\textbf{Tirocinio - Universit\`a degli studi di Trento, Dipartimento CIBIO} \hfill Marzo 2021 - Agosto 2021\\
Analisi di dati NGS di Trascrittomica \href{https://github.com/giacThePhantom/Tesi}{github.com/giacThePhantom/Tesi} \hfill \textit{Trento, Italia}\\
 \begin{itemize}
    \itemsep -3pt {}
     \item Implementazione di una pipeline per il processamento di dati RNA-seq su server remoto.
 \end{itemize}


\end{rSection}

%----------------------------------------------------------------------------------------
%	WORK EXPERIENCE SECTION
%----------------------------------------------------------------------------------------

\begin{rSection}{Attivit\`a di Mentoring e Formazione}
\vspace{-1.25em}

\item \textbf{Tutorato Didattico - Corso di informatica (Set. 2025 - Dic. 2025): }{
Supporto alla docenza per il corso di Informatica presso il dipartimento CIBIO, con focus sull'assistenza agli studenti nell'apprendimento della programmazione e della risoluzione di problemi algoritmici.
   }
\item \textbf{Co-supervisione di Tesi: }{Attività di tutoraggio tecnico-scientifico per due studenti durante lo svolgimento del loro lavoro di tesi sperimentale. Ho coordinato la pianificazione delle analisi e la formazione sull’utilizzo di strumenti bioinformatici, supervisionando il progresso dei progetti fino alla discussione finale.}
\end{rSection}

\begin{rSection}{Progetti Personali}
\vspace{-1.25em}

\item \textbf{Raccolta didattica open-source: }{
      Redazione integrale in \LaTeX del materiale didattico relativo ai corsi universitari da me seguiti, reso accessibile pubblicamente su: {\href{https://github.com/giacThePhantom}{github.com/giacThePhantom}}.
      L'iniziativa ha garantito ai colleghi l'accesso a risorse di alta qualit\`a, favorendo la nascita di una piccola community di contributori per il mantenimento e l'aggiornamento del progetto.
   }
\end{rSection}

%----------------------------------------------------------------------------------------

%----------------------------------------------------------------------------------------
% TECHINICAL STRENGTHS
%----------------------------------------------------------------------------------------
\begin{rSection}{SKILLS}

\begin{tabular}{ @{} >{\bfseries}l @{\hspace{6ex}} l }
Programmazione & Python (numpy, pandas, scikit-learn, pytorch), R, C/C++, Java, bash, Matlab, \\
 &Javascript, Sistemi Linux, Git, \LaTeX\\
Lingue & Fluente in Italiano e Inglese (C1)
\end{tabular}\\
\end{rSection}


\end{document}
